\begin{statementbox}
  \textbf{\textcolor{CStatement}{这条结论解决什么问题？}}
  已经画出$y=f(x)$，如果想让\textbf{整条图像升高或降低}，不用重新描点，直接改解析式就行。

  \medskip
  \textbf{\textcolor{CStatement}{使用条件}}
  \begin{itemize}[leftmargin=*,itemsep=0.25em,topsep=0.25em]
    \item 原函数是$y=f(x)$，图像\textbf{只上下移动}，不左右移动、不改变形状。
    \item 平移距离$b>0$；向上或向下由题目给定。原函数的定义域可以是任意集合，不要求$x$能取所有实数。
  \end{itemize}

  \textbf{\textcolor{CStatement}{核心结论：上加下减}}
  \[
    \boxed{\text{向上移 }b\text{ 个单位：}\ y=f(x)+b}
    \qquad
    \boxed{\text{向下移 }b\text{ 个单位：}\ y=f(x)-b}
  \]
  \textbf{为什么？}原来图像上有点$(x,f(x))$，上移后就是$(x,f(x)+b)$：\textbf{$x$不变，$y$加$b$}；下移同理减$b$。

  \begin{center}
  \begin{tikzpicture}[x=1.00cm,y=0.63cm,>=stealth]
    \draw[->,gray] (-2.20,0)--(2.40,0) node[right]{$x$};
    \draw[->,gray] (0,-1.55)--(0,5.45) node[above]{$y$};
    \draw[gray!75,thick,domain=-1.50:1.50,samples=60] plot (\x,{\x*\x});
    \draw[CStatement,very thick,domain=-1.50:1.50,samples=60] plot (\x,{\x*\x+2});
    \draw[blue!70!black,very thick,domain=-1.50:1.50,samples=60] plot (\x,{\x*\x-1});
    \draw[->,CStatement,thick] (0.20,0.25)--(0.20,2.25);
    \draw[->,blue!70!black,thick] (-0.25,-0.1)--(-0.25,-1.1);
    \fill (0,0) circle (1.4pt);
    \node[right,gray!80!black] at (1.9,1.65) {$y=x^2$};
    \node[right,CStatement] at (1.9,4.00) {$y=x^2+2$};
    \node[right,blue!70!black] at (1.9,-.5) {$y=x^2-1$};
    \node[left,CStatement] at (-2.3,2) {上移$2$};
    \node[left,blue!70!black] at (-2.3,-1) {下移$1$};
    \node[left] at (-.4,.2) {$O$};
  \end{tikzpicture}

  \end{center}
  \textbf{看图记忆：}图像上升$2$格，所有函数值都增加$2$；下降$1$格，所有函数值都减少$1$。

  \medskip
  \textbf{\textcolor{CStatement}{一种写法记住上下两种方向}}
  \[
    y=f(x)+t,\quad t\in\mathbb R.
  \]
  $t>0$上移$t$个单位；$t<0$下移$|t|$个单位；$t=0$不动。这里$t$是\textbf{带方向的变化量}，而移动的\textbf{距离始终非负}。
\end{statementbox}
